\chapter{Conclusions \& Future Scope}

Project Koschei successfully demonstrates the feasibility of deploying a memory-augmented Large Language Model as a real-time social deception agent within a multiplayer game environment. The system integrates a Unity-based first-person game client, a FastAPI backend, a ChromaDB vector memory store, and a locally hosted Ollama LLM into a cohesive architecture that enables the alien impostor to engage players through proximity-based chat, adapt its dialogue using stylometric mimicry, and maintain consistent deceptive narratives across multiple conversational rounds. The backend terminal results confirmed that the LLM pipeline correctly classified player messages, selected contextually appropriate deception strategies such as \texttt{seed\_doubt}, and generated responses within sub-second latency under normal load. The implemented gameplay systems --- including the dual-radius proximity grouping, conversation lifecycle management, NPC researcher interactions, and FSM-driven enemy and impostor behaviour --- collectively produced an immersive and psychologically tense multiplayer experience consistent with the project's design goals.

The project also establishes a meaningful contribution to the broader research areas of LLM-driven autonomous agents, real-time AI--human interaction, and social deception modelling. By replacing scripted impostor logic with a context-aware language model grounded in retrieved memory, Project Koschei moves beyond the fixed behavioural patterns of traditional social deduction games toward a system capable of open-ended, adaptive deception. The modular architecture ensures that individual components --- the game client, session layer, AI backend, and memory system --- can be independently developed, tested, and scaled, providing a robust foundation for continued research and feature development.

\vspace{9mm}
\noindent\textbf{\large Future Scope}
\vspace{3mm}

\par
Several avenues exist for extending Project Koschei beyond its current implementation. The LLM backbone could be upgraded to larger models such as \texttt{llama3.3:70b} or fine-tuned on curated social deception dialogue datasets to produce more nuanced and harder-to-detect impostor behaviour. A dynamic trust graph system could be introduced to quantify inter-player suspicion levels in real time, allowing the impostor to prioritise accusation deflection toward the most suspicious player and enabling richer emergent social dynamics. Streaming inference via WebSockets or server-sent events could replace the current request--response model to reduce perceived response latency and make NPC dialogue feel more natural and conversational. The game world could be expanded with additional map zones, more NPC researcher personalities with distinct LLM-driven dialogue profiles, and secondary alien entity types with varied deception strategies to increase replayability. Finally, integrating text-to-speech synthesis for impostor and NPC responses, combined with voice-based proximity chat for human players, would substantially enhance immersion and open new research directions in multimodal social deception detection.
